% Irodov Arihant PDF pp. 27--28, printed pp. 18--19.
% 111--117: Coriolis, rotating frames, Earth's rotation.
\exproblem{111}{18}
在地理纬度 $\varphi=60^\circ$ 处，步枪向正北方的
竖直标线水平射击，子弹初速为 $v=900\,\mathrm{m/s}$，
靶距 $s=1.0\,\mathrm{km}$。
忽略空气阻力，求子弹命中位置偏离标线的水平距离及方向。
\exsolution
取局部坐标 $\bm e_E$ 向东、$\bm e_N$ 向北、$\bm e_U$ 向上，
地球角速度矢量为
$\bm\Omega=\Omega(\cos\varphi\,\bm e_N+
\sin\varphi\,\bm e_U)$。
一阶近似忽略地球自转对弹道时间的影响，
子弹主要速度是 $\bm v=v\,\bm e_N$。
旋转参考系中的科里奥利加速度为
\[
\bm a_{\rm Cor}=-2\bm\Omega\times\bm v
=2\Omega v\sin\varphi\,\bm e_E.
\]
飞行时间 $t\simeq s/v$，故东向偏移
\[
\boxed{x_E\simeq\frac12(2\Omega v\sin\varphi)
\left(\frac sv\right)^2
=\frac{\Omega s^2\sin\varphi}{v}}.
\]
取 $\Omega=7.292\times10^{-5}\,\mathrm{rad/s}$，
得到 $\boxed{x_E\simeq0.070\,\mathrm m=7.0\,\mathrm{cm}}$，
方向\emph{向东}。
\excheck 该题只问相对竖直标线的水平偏差，
不是重力引起的落弹高度。
保留子弹落下造成的次阶东向修正约为
$O(\Omega gt^3)$，远小于主项。

\exproblem{112}{18--19}
水平转盘以恒角速度 $\omega=6.0\,\mathrm{rad/s}$
绕中心竖直轴旋转。
质量 $m=0.50\,\mathrm{kg}$ 的小物体沿转盘直径
以相对转盘恒速 $v'=0.50\,\mathrm{m/s}$ 运动。
当它距转轴 $r=0.30\,\mathrm m$ 时，
求转盘施加于小物体的合接触力。
\exsolution
在随盘旋转坐标系中，物体的相对加速度为零。
水平接触力需补偿离心力和科里奥利力：
\[
\bm F_{\rm horiz}
=-m\omega^2r\,\bm e_r+
2m\omega v'\,\bm e_\varphi ,
\]
其中方位符号随相对运动方向改变，两项恒垂直。
竖直方向接触支持力为 $N_z=mg$。因此盘对物体的
三维接触合力大小为
\[
\boxed{F=m\sqrt{g^2+\omega^4r^2+4\omega^2{v'}^2}}
\simeq7.89\,\mathrm N\approx8.0\,\mathrm N.
\]
\excheck 必须包括竖直法向支持力 $mg$；
只合成离心与科里奥利的水平项会低估题目的“盘施加的合力”。

\exproblem{113}{19}
一根光滑水平杆绕一端 $A$ 以恒定角速度
$\omega=2.00\,\mathrm{rad/s}$ 绕竖直轴转动。
质量 $m=0.50\,\mathrm{kg}$ 的滑套从 $A$ 处开始，
初始沿杆向外的相对速度 $v_0=1.00\,\mathrm{m/s}$。
当它距转轴 $r=0.50\,\mathrm m$ 时，
求随杆参考系中的科里奥利惯性力大小。
\exsolution
滑套的径向运动在旋转坐标系满足
$\ddot r=\omega^2r$；
因为杆光滑，杆没有沿自身的接触分力，
而离心力为 $m\omega^2r$。
把方程乘以 $2\dot r$ 并积分：
\[
\frac{\mathrm d}{\mathrm dt}
(\dot r^2-\omega^2r^2)=0,
\quad
\dot r^2=v_0^2+\omega^2r^2.
\]
科里奥利力垂直于杆，大小
$F_{\rm Cor}=2m\omega|\dot r|$，所以
\[
\boxed{F_{\rm Cor}=
2m\omega\sqrt{v_0^2+\omega^2r^2}
\simeq2.83\,\mathrm N}.
\]
\excheck 杆的实际接触力须提供相应横向加速度；
不能把“杆光滑”误解为其横向支持力为零。

\exproblem{114}{19}
一水平圆盘半径为 $R$，绕盘缘的竖直固定轴
以角速度 $\omega$ 转动。
质量为 $m$ 的粒子沿盘周以相对盘恒定速率运动。
在粒子离转轴最远时，随盘参考系中的惯性力合力恰为零。
求：（a）粒子的相对盘向心加速度；
（b）惯性力合力大小随其离转轴距离 $r$ 的关系。
\exsolution
取转轴为原点，盘心在 $(R,0)$。
盘周粒子相对盘的位置可写为
$\bm r=R(1+\cos\theta,\sin\theta)$，
距转轴 $r=2R|\cos(\theta/2)|$。
在随盘参考系中，
\[
\bm F_{\rm in}=
m\omega^2\bm r-2m\bm\Omega\times\bm v'.
\]
最远点 $\theta=0$、$r=2R$ 时，惯性力要抵消，
必须沿盘周\emph{逆于盘转向}运动，且满足
$2m\omega v'=m\omega^2(2R)$，即
$v'=\omega R$。因此相对盘沿半径 $R$ 圆周的向心加速度为
\[
\boxed{a'_n=\frac{{v'}^2}{R}=\omega^2R}.
\]
将相对盘的速度取为
$\bm v'=-\omega R(-\sin\theta,\cos\theta)$，
代回惯性力表达式，整理为
$\bm F_{\rm in}
=m\omega^2R(1-\cos\theta,-\sin\theta)$，其大小
\[
\boxed{|\bm F_{\rm in}(r)|=
m\omega^2\sqrt{4R^2-r^2}}.
\]
\excheck 当 $r=2R$ 合力为零；
当 $r=0$，其大小为 $2m\omega^2R$，
与盘心移动及科里奥利项的矢量合成一致。

\exproblem{115}{19}
质量为 $m=0.30\,\mathrm{kg}$ 的小物块从光滑球体
顶部由静止开始滑下，球半径 $R=1.00\,\mathrm m$，
球体以恒角速度 $\omega=6.0\,\mathrm{rad/s}$
绕经过球心的竖直轴旋转。
在随球参考系中，求离面瞬间离心惯性力与科里奥利惯性力大小。
\exsolution
因为球面光滑且轴对称，物块在惯性系中不受绕竖直轴的切向力；
从顶点释放时方位角动量为零，故物块沿固定的竖直子午面滑动，
其离面条件与不转动球的第87题相同：
\[
\cos\theta=\frac23,\quad
\sin\theta=\frac{\sqrt5}{3},\quad
v^2=\frac{2gR}{3}.
\]
离轴水平距离为
$r=R\sin\theta=R\sqrt5/3$，所以离心力
\[
\boxed{F_{\rm cf}=m\omega^2r
=\frac{\sqrt5}{3}m\omega^2R\simeq8.05\,\mathrm N}.
\]
物块在旋转系中的相对速度
$\bm v'=\bm v-\bm\Omega\times\bm r$。
其水平分量有径向 $v\cos\theta$ 与方位向
$-\omega r$ 两个互相垂直的分量，
所以
\[
F_{\rm Cor}=2m\omega\sqrt{v^2\cos^2\theta+\omega^2r^2}
=\boxed{\frac{2m\omega R}{3}
\sqrt{5\omega^2+\frac{8g}{3R}}}
\simeq17.2\,\mathrm N .
\]
\excheck 科里奥利力要用\emph{相对旋转球体}的速度，
不能把惯性系中沿子午线的 $v$ 直接代入 $2m\omega v$。

\exproblem{116}{19}
质量为 $m=2000$ 吨的列车位于北纬 $\varphi=60^\circ$。
求：（a）它以 $v=54\,\mathrm{km/h}$ 沿子午线行驶时，
对钢轨的横向作用力大小和方向；
（b）列车沿哪个方向、以多大速率运动，
可使地球参考系中的离心力与科里奥利惯性力之和为零。
\exsolution
（a）以 $m=2.0\times10^6\,\mathrm{kg}$、
$v=15\,\mathrm{m/s}$ 和
$\Omega=7.292\times10^{-5}\,\mathrm{rad/s}$ 计算，
沿子午线速度产生的科里奥利横向加速度大小是
$2\Omega v\sin\varphi$。列车对轨道的横向作用力
\[
\boxed{F_\perp=2m\Omega v\sin\varphi
\simeq3.79\,\mathrm{kN}}.
\]
在北半球，无论列车向北或向南行驶，科里奥利惯性力
相对于行车方向都向右，因此列车压\emph{右侧钢轨}。
（b）地球自转的离心惯性力矢量与赤道半径的投影成正比，
量级为 $m\Omega^2R_\oplus\cos\varphi$，方向背离地轴。
若列车以速率 $u$ 沿纬线向西行驶，科里奥利惯性力
与该离心力矢量反向；令两者大小相等：
\[
2m\Omega u=m\Omega^2R_\oplus\cos\varphi.
\]
故
\[
\boxed{u=\frac12\Omega R_\oplus\cos\varphi
\simeq116\,\mathrm{m/s}\simeq418\,\mathrm{km/h}},
\]
方向\emph{沿纬线由东向西}。
\excheck （b）比较的是完整的离心和科里奥利惯性力矢量，
不是沿轨道运动时的牵引力。用 $R_\oplus\simeq6371\,\mathrm{km}$
估算；真实地球有效重力与形状修正未包含在理想模型中。

\exproblem{117}{19}
在赤道，某物体相对地面初始静止，
从高度 $h=500\,\mathrm m$ 自由下落，忽略空气阻力。
求它落地时相对铅垂线的水平偏移及方向。
\exsolution
采用当地向东、向北、向上坐标。
赤道处 $\bm\Omega=\Omega\bm e_N$。
在忽略自转的零阶运动中，
竖直速度 $v_z=-gt$，所以
科里奥利加速度的东向分量为
\[
a_E=-2(\bm\Omega\times\bm v)_E=2\Omega gt.
\]
从初始水平速度零积分：
$v_E=\Omega gt^2$、$x_E=\Omega gt^3/3$。
自由落体时间 $t_f=\sqrt{2h/g}$，代入
\[
\boxed{x_E\simeq\frac{\Omega g}{3}
\left(\frac{2h}{g}\right)^{3/2}
=\frac23\Omega h\sqrt{\frac{2h}{g}}
\simeq0.245\,\mathrm m},
\]
方向\emph{向东}。
\excheck 只保留地球自转角速度的一阶效应；
若忽略科里奥利力会错误给出零水平偏差。
